Al fine di trovare il punto di equilibrio del sistema, è necessario considerare nulle tutte le derivate del sistema.\\
Poniamo quindi le derivate dell'equazione \eqref{eq:model} pari a zero e otteniamo il seguente risultato:
\begin{displaymath}
\left\{
\begin{array}{l}
0=\bar F \cos (\bar\vartheta)-\bar F_L \operatorname{sin}(\bar\vartheta)-\bar B \\
0=0-k_{\vartheta} \operatorname{sin}(\bar\vartheta)-\bar{F} b \\
y=\bar{z}-b \operatorname{sin}(\bar\vartheta)
\end{array}
\right.
\end{displaymath}
Svolgendo i calcoli otteniamo:
\begin{displaymath}
\left\{
\begin{array}{l}
\bar F_L= \frac{\bar F \cos(\bar\vartheta)-\bar B}{\sin(\bar\vartheta)} \\
\bar F = -\frac{k_\vartheta \sin(\bar\vartheta)}{b} \\
\bar z = y + b\sin(\bar\vartheta)
\end{array}
\right.
\end{displaymath}
Sostituendo $\bar F$ in $\bar F_L$:
\begin{displaymath}
\left\{
\begin{array}{l}
\bar F_L= \frac{(-\frac{k_\vartheta \sin(\bar\vartheta)}{b}) \cos(\bar\vartheta)-\bar B}{\sin(\bar\vartheta)} \\
\bar F = -\frac{k_\vartheta \sin(\bar\vartheta)}{b} \\
\bar z = y + b\sin(\bar\vartheta) \\
\end{array}
\right.
\end{displaymath}
Sostituendo il valore $\bar\vartheta = \frac{\pi}{6}$:
\begin{displaymath}
\left\{
\begin{array}{l}
\bar F_L= 2(-\frac{\sqrt3}{4}\frac{k_\vartheta}{b}-\bar B)\\
\bar F = -\frac{k_\vartheta}{2b} \\
\bar z = y + \frac{b}{2} \\
\end{array}
\right.
\end{displaymath}